% ID: FIS-TER-GAS-003
% Titolo: Calore fornito in un'espansione isoterma
% Disciplina: Fisica
% Area: Termologia
% Argomento: Gas perfetti
% Sottoargomento: Espansione isoterma
% Classe: Quarta liceo scientifico
% Difficolta: 4
% Tipo: problema_numerico
% Risultato: Q = 4,2 * 10^2 J
% Tempo_stimato: 10 min
% Competenze: applicazione del primo principio, calcolo del lavoro isotermo, interpretazione energetica
% Prerequisiti: gas perfetti, trasformazione isoterma, logaritmi
% Tag: termologia, gas_perfetti, isoterma, calore, lavoro
% Fonte: verifica_fisica_gas_perfetti_anonimizzata
% Autore: Riccardo Carli
% Licenza: CC-BY-SA-4.0

\begin{esercizio}
Un cilindro contiene \(0{,}30\,\mathrm{mol}\) di gas perfetto alla temperatura
di \(300\,\mathrm{K}\). Il gas si espande isotermicamente da un volume iniziale
di \(0{,}20\,\mathrm{m^3}\) a un volume finale di \(0{,}35\,\mathrm{m^3}\).

Quanto calore deve essere fornito al gas per mantenere costante la temperatura?
\end{esercizio}

\begin{soluzione}
In una trasformazione isoterma di un gas perfetto la variazione di energia
interna e nulla:
\[
\Delta U=0.
\]
Dal primo principio segue quindi che il calore assorbito e uguale al lavoro
compiuto dal gas:
\[
Q=L.
\]

Per un'espansione isoterma reversibile:
\[
L=nRT\ln\frac{V_f}{V_i}.
\]
Sostituendo i dati:
\[
Q=0{,}30\cdot 8{,}31\cdot 300
\ln\frac{0{,}35}{0{,}20}.
\]

\[
Q\simeq 4{,}2\cdot 10^2\,\mathrm{J}.
\]

Il calore deve essere fornito al gas, quindi il suo segno e positivo.
\end{soluzione}
